%% kapa: exam=1 round=26 year=2015 period=1 form=A subject=civil_law subject_ko=민법 seq=22 number=22
%% source: Q-Net 감정평가사 기출문제 (cst00309) artlSeq=5138001 file=extracted/1008/5138001/제26회_감정평가사 1차_기출/제26회_감정평가사 1차_기출/1교시_감평_A형.pdf
%% conversion: kordoc 4.24.51 -> tools/gen1.py (auto); review=proofread-sonnet; flags=none
\begin{question}{22}
물권적 청구권에 관한 설명으로 옳지 않은 것은? (다툼이 있으면 판례에 따름)
\begin{choices}
  \item 甲의 물건을 乙이 불법 점유하는 경우 甲은 丙에게 그 소유권을 양도하면서 乙에 대한 소유물반환 청구권을 자신에게 유보할 수 없다.
  \item 소유자는 현재 점유하고 있지 않은 자를 상대로 소유물의 반환을 청구할 수 없다.
  \item 물권적 청구권은 점유권과 소유권 이외의 물권에 대하여도 인정된다.
  \item 소유권에 기한 물권적 청구권은 소멸시효에 걸리지 않는다.
  \item 간접점유자는 직접점유자가 점유의 침탈을 당한 때에는 그 물건의 반환을 청구할 수 없다.
\end{choices}
\end{question}
\answer{5}
